\section{Additional Related Work, Fair Comparison, and Limitations}
\label{sec:related-limitations}

\subsection{Related perspectives}

\textbf{Support-aware policy improvement.}
Conservative offline policy-improvement methods limit updates in regions where
the behavior data provide weak support.  SPIBB constrains uncertain decisions
to the baseline policy~\cite{spibb2019}, while BEAR regularizes a learned policy
toward the behavior distribution~\cite{bear2019}.  PC-MTS shares the support
intuition but applies it before optimization: it filters trajectory-level
supervision using samples from the frozen learner, with command-conditional
calibration and a local feasibility test.  It is not a general monotonic-
improvement guarantee; its role is to avoid moving the imitation policy into a
region where downstream rollout optimization has poor local support.

\textbf{Multi-objective and Pareto optimization.}
Multi-objective reinforcement learning studies policies under vector-valued
returns and alternative preference models~\cite{roijers2013}.  AMPT does not
seek a global set of Pareto-optimal policies.  PC-MTS uses a scene-local Pareto
filter for supervision, and FF-PGRPO uses a hierarchy: hard feasibility and
coherent-reference guards are evaluated before a rollout can receive positive
Pareto credit.  This order is essential when an aggregate score can compensate
for a protected regression.

\textbf{Advantage-weighted behavior learning.}
Advantage-weighted regression turns estimated advantages into weights for
supervised policy updates~\cite{awr2021}.  APR is related at the loss level,
but its weights are attached only after a candidate is compared with the
current single-trajectory prediction, locally interpolated, and re-evaluated.
Teacher admission and policy retention are therefore as important as the
weighting function itself.

\textbf{Model merging.}
Task arithmetic combines model displacements relative to a shared
anchor~\cite{taskarithmetic2023}, and TIES resolves redundancy and sign
conflicts before merging~\cite{ties2023}.  AMPT may use an anchor-constrained
variant to consolidate APR branches, but merging is optional and is separated
from teacher selection.  Passing parameter-space norm and sign checks does not
imply behavioral safety; every merged checkpoint must still undergo the same
single-trajectory evaluator audit.

\subsection{Fair-comparison boundary}
\label{sec:fair-comparison}

AMPT uses one predicted trajectory at inference.  Policy rollout banks,
candidate pools, coherent references, group sampling, specialist policies, and
the offline evaluator are training-only; there is no test-time scorer,
best-of-$N$ selection, or candidate reranking.  This contract is present in the
paper and is independently supported by the v1 submission provenance.

The released v1 scene export contains 12,138 successfully scored scenes from
12,146 predictions; the eight unscored predictions are disclosed rather than
imputed.  The v2 export contains 12,146 successful scene scores.  Its EC
aggregation follows the official evaluator, although a direct EC value is
present for 10,040 rows.  Scene-bootstrap intervals resample these scored
scenes and are explicitly labeled as scene uncertainty; they are not a
substitute for independent training runs.

The baseline rows in the main tables are treated as \emph{reported} unless a
local scene export and matching protocol are available.  AMPT's v1/v2 rows are
reproduced from scene-level outputs, but protocol matching is not established
for every external baseline.  Moreover, the audited v1 provenance records use
of navtest for checkpoint selection.  We therefore avoid a strict held-out
state-of-the-art claim and report protocol fields---backbone, sensors, training
data, inference candidates, and test-time scoring---alongside each comparison.

The 658-scene result uses the recovery checkpoint and checkpoint pair specified
by the main-paper analysis.  It is a fixed-subset diagnostic, not a claim that
every later checkpoint must induce the same transitions.  Reporting the full
$2\times2$ matrix in Eq.~\ref{eq:hard658-transition} makes the gross and net
notions of recovery explicit.

\subsection{Limitations}
\label{sec:limitations}

\textbf{Finite support estimates.}
Compatibility is estimated from a finite rollout bank.  Sparse commands or
multi-modal scenes can make a genuinely useful target appear distant, and the
calibrated neighborhood is only as representative as the frozen policy
samples.  Increasing bank size improves coverage but raises offline inference
and storage cost.

\textbf{Evaluator bias and non-reactive evaluation.}
All three stages use an offline NAVSIM evaluator for filtering, credit, or
teacher validation.  Systematic evaluator errors can therefore be amplified
across stages.  NAVSIM's non-reactive agents also cannot reveal how other road
users would respond to the learned trajectory; feasibility under this scorer
is not a closed-loop safety guarantee.

\textbf{Candidate-pool ceiling.}
APR can select and localize only behaviors represented in its multi-source
pool.  If no source proposes a valid maneuver, rebuilding the pool changes no
target.  The persistent case in Table~\ref{tab:hard658-cases} establishes
persistence under two audited checkpoints, but cannot diagnose this pool
ceiling without a per-scene candidate manifest.

\textbf{Locality can reject long-horizon value.}
The policy-compatibility filter intentionally favors learnable local targets.
It may reject a distant behavior that is initially difficult to imitate but
would be valuable after a longer curriculum.  Policy-dependent teacher-value
and cross-policy selection experiments are therefore important follow-ups.

\textbf{Benchmark-specific hierarchy.}
The ordering of NC/DAC feasibility, DDC/TLC guards, progress floors, and
continuous Pareto objectives reflects NAVSIM metrics.  Transferring AMPT to a
new benchmark requires revalidating the hierarchy and calibration thresholds,
not merely renaming metric fields.

\textbf{No behavioral semantics for parameter merging.}
Task-vector masks, direction agreement, and norm constraints bound parameter
movement but do not provide a semantic or safety guarantee.  If its independent
gain is not stable across seeds and step-matched controls, merging should
remain an implementation detail.

\textbf{Statistical scope.}
The supervision comparison, stage ablation, and APR round table are available
as reported point estimates rather than three-seed aggregates.  Small APR
increments must not be described as stable until the configured multi-seed
controls are run.  The paired 658-scene interval establishes a robust scene-
level net transition for one checkpoint pair, but does not measure training-
seed variation.

These limitations narrow rather than alter the paper's central claim.  The
currently traceable evidence establishes the supervision--optimization
reversal, a large feasibility-oriented stage gain, monotonic reported APR
rounds with preserved displayed safety components, and positive net recovery
on a fixed hard set.  Candidate-count matching, rollout-level credit traces,
step-matched APR controls, and multi-seed reruns are the remaining experiments
needed for stronger causal and stability claims.
